% ============================================================================
% SIMULADO E REVISÃO INTEGRADA DO MÓDULO 1 (CAPÍTULOS 1 A 4)
% Componente: CHF662 — Análise Ambiental (DCHF/UEFS)
% Vínculo no Plano CHF662: Revisão EO1–EO3 e Preparação para Prova Teórica
% ============================================================================

\clearpage
\phantomsection
\addcontentsline{toc}{chapter}{Simulado e Revisão Integrada: Módulo 1}
\chapter*{Simulado e Revisão Integrada: Módulo 1}
\markboth{SIMULADO DO MÓDULO 1}{SIMULADO DO MÓDULO 1}

\noindent\tagEmenta\ \tagEO{1}\ \tagEO{2}\ \tagEO{3}

\vspace{6pt}

\begin{boxConceito}[Instruções de Aplicação e Alinhamento Avaliativo]
Este simulado integra os conteúdos dos \textbf{Capítulos 1 a 4} (Semanas 1 a 4 do Plano CHF662), preparando diretamente para a consolidação dos \textbf{Estudos Orientados (EO1, EO2 e EO3)} e para a \textbf{Avaliação Teórica Individual} (Semana 10).
\begin{itemize}[leftmargin=1.5em, itemsep=2pt]
  \item \textbf{Tempo sugerido:} 90 minutos (50 minutos para as 8 questões objetivas + 40 minutos para as 2 questões dissertativas).
  \item \textbf{Modo de estudo:} resolva todas as questões sem consulta. Ao concluir, utilize o \textbf{Gabarito Integralmente Comentado}, o \textbf{Espelho Dissertativo} e a \textbf{Tabela de Diagnóstico de Revisão} na página seguinte para identificar exatamente quais capítulos revisar.
  \item \textbf{Pontuação simulada (escala $0$ a $10,0$):} Questões 1 a 8 valem $0,75$ ponto cada ($6,0$ pontos); Questões 9 e 10 valem $2,0$ pontos cada ($4,0$ pontos).
\end{itemize}
\end{boxConceito}

% ============================================================================
% PARTE I — QUESTÕES OBJETIVAS CONTEXTUALIZADAS (1 A 8)
% ============================================================================
\section*{Parte I: Questões Objetivas (Estilo ENADE e Concursos Ambientais)}

\begin{enumerate}[leftmargin=1.8em, itemsep=6pt]

  \item \textbf{(Cap. 1, Questão autoral, estilo ENADE / Geografia)} Ao discutir a crise ambiental contemporânea, Enrique Leff (2001) sustenta que a degradação generalizada dos solos, a escassez hídrica e a perda de biodiversidade não podem ser solucionadas apenas por ajustes tecnológicos pontuais. De acordo com essa perspectiva epistemológica, a crise ambiental contemporânea caracteriza-se fundamentalmente como:
  \begin{alternativas}
    \item Um ciclo climático puramente geológico, independente das formas de organização econômica e técnica da humanidade.
    \item Uma crise civilizatória e de racionalidade científica e econômica, produzida por um modelo que fragmentou o saber e tratou a natureza como recurso ilimitado.
    \item Um desequilíbrio restrito a países de clima semiárido, solucionável apenas pela importação de tecnologia de irrigação intensiva.
    \item Uma consequência inevitável do crescimento demográfico absoluto, desvinculada dos padrões de consumo e da desigualdade territorial.
    \item Um problema estritamente jurídico, cuja resolução depende apenas da ampliação do valor das multas administrativas sem mudança produtiva.
  \end{alternativas}
  \espacoResolucao[1.2cm]

  \item \textbf{(Caps. 1 e 2, Questão autoral, estilo ENADE / Geografia)} Carlos Walter Porto-Gonçalves (2006) analisa como o pensamento moderno consolidou a ideia de ``natureza'' como algo externo à humanidade, passível de domínio técnico irrestrito. Sobre as consequências dessa concepção moderna na relação sociedade--natureza, assinale a alternativa correta:
  \begin{alternativas}
    \item A modernidade integrou plenamente os saberes tradicionais e os ciclos ecológicos locais ao ritmo de acumulação industrial.
    \item A noção de natureza é universal e imutável em todas as culturas, não sofrendo influência histórica, política ou simbólica.
    \item A separação sujeito--objeto legitimou a apropriação mercantil dos recursos naturais e invisibilizou modos de vida de populações que se concebem como parte dos ecossistemas.
    \item A Geografia moderna surgiu para rejeitar completamente o estudo da materialidade natural em favor da filosofia abstrata.
    \item A concepção moderna eliminou os conflitos socioambientais ao garantir distribuição homogênea de água e terra no espaço rural brasileiro.
  \end{alternativas}
  \espacoResolucao[1.2cm]

  \item \textbf{(Cap. 2, Questão autoral, estilo Concurso Analista Ambiental / INEMA)} Em uma região do Semiárido baiano, a instalação de um empreendimento minerário de grande porte alterou o acesso de comunidades tradicionais de fundo de pasto a aguadas comunitárias e áreas de uso coletivo da Caatinga. A empresa classificou a reação comunitária como ``entrave cultural ao progresso''. À luz do debate sobre conflitos socioambientais (Porto-Gonçalves, 2006; Mendonça, 2001), esse episódio expressa:
  \begin{alternativas}
    \item Um conflito socioambiental territorial, no qual distintos grupos sociais disputam a apropriação material, simbólica e decisória da base ecológica local.
    \item Uma divergência puramente técnica que deve ser resolvida sem participação social, bastando apresentar mapas geológicos do minério.
    \item Um problema restrito à Geografia Física, pois envolve apenas a recarga hídrica subterrânea das aguadas.
    \item Uma comprovação do determinismo ambiental, visto que o clima semiárido impede qualquer convivência produtiva sustentável.
    \item Um impacto exclusivamente natural, sem relação com modelos de desenvolvimento ou desigualdade de poder.
  \end{alternativas}
  \espacoResolucao[1.2cm]

  \item \textbf{(Cap. 3, Questão autoral, estilo ENADE / Geografia)} Sobre a proposta da \textbf{Geografia Socioambiental} formulada por Francisco Mendonça (2001) e sua posição frente à história do pensamento geográfico moderno, avalie as afirmações a seguir:
  \begin{enumerate}[label=\Roman*., leftmargin=2.0em, itemsep=2pt]
    \item A proposta surge como resposta crítica à dicotomia entre uma Geografia Física excessivamente naturalista e uma Geografia Humana que secundarizou a materialidade ecológica.
    \item O prefixo ``socio'' indica que os problemas ambientais resultam da interação indissociável entre processos biofísicos e organização social, política e econômica.
    \item A abordagem socioambiental recomenda que o relatório de diagnóstico seja dividido em anexos isolados de meio físico e meio humano sem cruzamento causal, para preservar a neutralidade de cada ramo.
  \end{enumerate}
  É correto o que se afirma apenas em:
  \begin{alternativas}
    \item I.
    \item II.
    \item I e II.
    \item II e III.
    \item I, II e III.
  \end{alternativas}
  \espacoResolucao[1.2cm]

  \item \textbf{(Cap. 3, Questão autoral, estilo Concurso Geógrafo / Planejamento Ambiental)} No diagnóstico integrado de uma microbacia urbana com lagoas assoreadas em Feira de Santana (BA), uma equipe aplicou o \textbf{Sistema Tripolar GTP} de Georges Bertrand (2004). Assinale a alternativa que associa corretamente um procedimento de campo ao respectivo polo do sistema:
  \begin{alternativas}
    \item Levantamento de escrituras, legislação de APP e disputas imobiliárias $\rightarrow$ Polo Geossistema (G).
    \item Medição da taxa de infiltração do solo, balanço hídrico e transporte de sedimentos sob influência da impermeabilização urbana $\rightarrow$ Polo Geossistema (G).
    \item Entrevistas sobre memória afetiva das lagoas e valorização simbólica do espelho d'água pelos moradores antigos $\rightarrow$ Polo Território (T).
    \item Mapeamento de declividade e granulometria dos sedimentos depositados no fundo da lagoa $\rightarrow$ Polo Paisagem (P), por tratar-se de dinâmica estritamente cultural.
    \item Fiscalização de loteamentos irregulares pela secretaria municipal de meio ambiente $\rightarrow$ Polo Paisagem (P), por envolver a estética dos muros urbanos.
  \end{alternativas}
  \espacoResolucao[1.2cm]

  \item \textbf{(Cap. 4, Questão autoral, estilo Concurso Analista Ambiental / IBAMA)} Durante a implantação de uma linha de transmissão elétrica atravessando vertentes íngremes e cursos d'água, a supressão da mata ciliar somou-se à compactação do solo por maquinário pesado. Essa combinação intensificou o escoamento concentrado e provocou ravinas muito mais profundas do que ocorreria se apenas um dos dois processos atuasse isoladamente. Conforme Sánchez (2013) e a Resolução CONAMA nº 01/1986, é correto afirmar que:
  \begin{alternativas}
    \item A formação de ravinas é o aspecto ambiental, enquanto a supressão da mata ciliar e a compactação são os impactos ambientais.
    \item O agravamento desproporcional resultante da interação entre dois processos distintos exemplifica o caráter \textbf{sinérgico} do impacto ambiental.
    \item Por não envolver lançamento de produto químico tóxico, a erosão em ravinas não constitui degradação ambiental nem impacto ambiental.
    \item A obra produz apenas impactos permanentes e irreversíveis, não sendo possível adotar medidas mitigadoras ou de recuperação.
    \item O impacto descrito classifica-se como benéfico e indireto, por ampliar a visibilidade dos afloramentos rochosos na bacia.
  \end{alternativas}
  \espacoResolucao[1.2cm]

  \item \textbf{(Cap. 4, Questão autoral, estilo Concurso Defesa Civil / Gestão de Riscos)} Um município baiano aplicou a matriz conceitual de risco ambiental $R = P \times V$ (escala ordinal de $1$ a $4$, com classes: $1$--$3$ Baixo; $4$--$6$ Médio; $8$--$9$ Alto; $12$--$16$ Crítico) a quatro setores de encosta e fundo de vale, obtendo os seguintes dados:
  \begin{center}
  \renewcommand{\arraystretch}{1.15}
  \small
  \begin{tabular}{@{}lcc@{}}
    \toprule
    \rowcolor{uefsNavy}
    \textcolor{white}{\bfseries Setor Avaliado} & \textcolor{white}{\bfseries Perigo ($P$)} & \textcolor{white}{\bfseries Vulnerabilidade ($V$)} \\
    \midrule
    Setor I (encosta íngreme ocupada sem drenagem) & $4$ & $4$ \\
    \rowcolor{bgGray}
    Setor II (margem de córrego com ocupação consolidada precária) & $3$ & $3$ \\
    Setor III (encosta muito íngreme com vegetação preservada e acesso restrito) & $4$ & $1$ \\
    \rowcolor{bgGray}
    Setor IV (área de baixa declividade com moradias de vulnerabilidade alta) & $2$ & $3$ \\
    \bottomrule
  \end{tabular}
  \end{center}
  Com base no cálculo exato de $R = P \times V$, assinale a alternativa correta:
  \begin{alternativas}
    \item Os Setores I e III pertencem ambos à classe Crítica, pois compartilham $P = 4$.
    \item O Setor II ($R = 9$, Alto) apresenta risco inferior ao do Setor IV ($R = 6$, Médio).
    \item O Setor I apresenta risco Crítico ($R = 16$), o Setor II apresenta risco Alto ($R = 9$), e os Setores III ($R = 4$) e IV ($R = 6$) enquadram-se na classe de risco Médio.
    \item A ordem decrescente de prioridade técnica de intervenção é: Setor I $>$ Setor III $>$ Setor II $>$ Setor IV.
    \item O Setor III deve ser classificado como Risco Baixo ($R = 4$) porque sua vulnerabilidade é igual a $1$.
  \end{alternativas}
  \espacoResolucao[1.5cm]

  \item \textbf{(Cap. 4, Questão autoral, estilo Concurso Perito Ambiental / IBAMA--INEMA)} Uma antiga área de mineração de cascalho em encosta do Semiárido sofreu decapagem completa do solo e rebaixamento topográfico de $6\text{ m}$, apresentando resiliência natural muito baixa. O órgão ambiental exige um PRAD que elimine a erosão ativa, estabilize os taludes e permita o aproveitamento seguro do terreno como viveiro comunitário de mudas nativas e área de captação de água pluvial, reconhecendo a impossibilidade física de reconstituir o relevo e a vegetação originais. O conceito operacional que define corretamente essa estratégia é:
  \begin{alternativas}
    \item Restauração ecológica integral.
    \item Poluição controlada.
    \item Reabilitação ambiental (enquadrada no conceito amplo de recuperação).
    \item Fragilidade potencial nula.
    \item Remediação química de metais pesados.
  \end{alternativas}
  \espacoResolucao[1.2cm]

\end{enumerate}

% ============================================================================
% PARTE II — QUESTÕES DISSERTATIVAS PREPARATÓRIAS (9 E 10)
% ============================================================================
\section*{Parte II: Questões Dissertativas (Preparação para EO1--EO3 e Prova Teórica)}

\begin{enumerate}[leftmargin=1.8em, itemsep=6pt, start=9]

  \item \textbf{(Integração Caps. 1, 2 e 3, Preparatória para EO1, EO2 e Prova Teórica)} Em uma microbacia do Semiárido baiano, a expansão de atividades intensivas sem planejamento desencadeou conflitos pelo uso da água, desmatamento de cabeceiras e perda de referência cultural dos moradores com o riacho local. Redija um texto técnico articulado (15 a 22 linhas) explicando:
  \begin{enumerate}[label=(\alph*), leftmargin=1.8em, itemsep=1.5pt]
    \item Como a crise ambiental nesse caso reflete a ruptura moderna entre sociedade e natureza (Leff, 2001; Porto-Gonçalves, 2006);
    \item Por que a Geografia Socioambiental (Mendonça, 2001) supera os limites das análises puramente naturalistas ou exclusivamente sociais;
    \item Como o \textbf{Sistema Tripolar GTP} de Georges Bertrand (2004) organizaria o diagnóstico integrado dessa microbacia nos polos Geossistema (G), Território (T) e Paisagem (P).
  \end{enumerate}
  \espacoResolucao[4.4cm]

  \item \textbf{(Integração Caps. 3 e 4, Preparatória para EO3 e Prova Teórica)} Durante a abertura de uma estrada de acesso em vertente declivosa da Caatinga sem sistema de drenagem, o solo exposto sofreu ravinamento severo e sedimentos foram carreados para um açude comunitário situado na base da encosta, onde vivem famílias dependentes da água armazenada. Com base em Sánchez (2013) e Tauk-Tornisielo (1996):
  \begin{enumerate}[label=(\alph*), leftmargin=1.8em, itemsep=1.5pt]
    \item Decomponha a cadeia causal identificando \textbf{atividade}, \textbf{aspecto ambiental}, \textbf{impacto direto} e \textbf{impacto indireto};
    \item Esclareça se o assoreamento do açude por sedimentos erodidos configura degradação ambiental, poluição ou ambos, justificando com base na PNMA;
    \item Considere que, no Cenário 1 (sem drenagem nem contenção), a área de jusante apresenta Perigo Crítico ($P_1 = 4$) e Vulnerabilidade Alta ($V_1 = 3$). No Cenário 2 (após implantação de drenagem, paliçadas retentoras de sedimentos e proteção das margens do açude), o perigo cai para Médio ($P_2 = 2$) e a vulnerabilidade para Média ($V_2 = 2$). Calcule $R_1$, $R_2$, suas respectivas classes na matriz $4 \times 4$ e a variação percentual $\Delta R\%$;
    \item Indique se a recomposição das margens do açude (APP) deve priorizar \textbf{restauração} ou \textbf{reabilitação}, justificando sua resposta.
  \end{enumerate}
  \espacoResolucao[4.4cm]

\end{enumerate}

% ============================================================================
% GABARITO INTEGRALMENTE COMENTADO, ESPELHO E TABELA DE DIAGNÓSTICO
% ============================================================================
\clearpage
\section*{Parte III: Gabarito Integralmente Comentado (Questões 1 a 8)}

\begin{boxGabarito}[Análise Comentada das Questões Objetivas]
\begin{enumerate}[leftmargin=1.8em, itemsep=4pt]
  \item \textbf{Gabarito: (B).} Para Enrique Leff (2001), a crise ambiental é uma crise civilizatória e de racionalidade (epistemológica e econômica) derivada da fragmentação do conhecimento e da lógica de acumulação que ignora limites ecológicos. \textit{Distratores:} (A) naturaliza a crise; (C) reduz a questão ao clima semiárido; (D) reproduz neomalthusianismo simplista; (E) reduz o problema a multas administrativas.
  \item \textbf{Gabarito: (C).} Porto-Gonçalves (2006) demonstra que a separação moderna sujeito--objeto transformou a natureza em objeto externo e mercadoria, deslegitimando saberes e territorialidades de povos integrados aos ecossistemas. \textit{Distratores:} (A) e (E) invertem o processo histórico; (B) naturaliza uma categoria histórica; (D) distorce a história da Geografia.
  \item \textbf{Gabarito: (A).} A disputa por aguadas e terras de uso coletivo envolvendo empresas e comunidades tradicionais caracteriza um conflito socioambiental territorial (apropriação material, simbólica e poder decisório). \textit{Distratores:} (B), (C) e (E) despolitizam e fragmentam o fenômeno; (D) ressuscita o determinismo ambiental.
  \item \textbf{Gabarito: (C).} As afirmações I e II reproduzem fielmente os fundamentos da Geografia Socioambiental em Mendonça (2001). A afirmação III é falsa porque anexos isolados sem cruzamento causal caracterizam a multidisciplinaridade justaposta que a Geografia Socioambiental busca superar.
  \item \textbf{Gabarito: (B).} Taxa de infiltração, balanço hídrico e transporte de sedimentos (incluindo o efeito físico da impermeabilização urbana) pertencem ao polo \textbf{Geossistema (G)}. \textit{Distratores:} (A) descreve Território, não G; (C) descreve Paisagem, não T; (D) descreve Geossistema, não P; (E) descreve Território, não P.
  \item \textbf{Gabarito: (B).} Quando dois processos distintos (supressão de mata ciliar + compactação por máquinas) interagem gerando efeito maior do que a soma isolada das partes, o impacto é \textbf{sinérgico}. \textit{Distratores:} (A) inverte aspecto e impacto; (C) confunde degradação com poluição química; (D) ignora a mitigação/recuperação; (E) classifica erroneamente como benéfico.
  \item \textbf{Gabarito: (C).} Cálculos passo a passo na matriz $R = P \times V$:
  \begin{itemize}[leftmargin=1.4em, itemsep=1pt]
    \item Setor I: $R_{\text{I}} = 4 \times 4 = 16$ (\textbf{Crítico});
    \item Setor II: $R_{\text{II}} = 3 \times 3 = 9$ (\textbf{Alto});
    \item Setor III: $R_{\text{III}} = 4 \times 1 = 4$ (\textbf{Médio});
    \item Setor IV: $R_{\text{IV}} = 2 \times 3 = 6$ (\textbf{Médio}).
  \end{itemize}
  Ordem decrescente real: $\text{I} (16) > \text{II} (9) > \text{IV} (6) > \text{III} (4)$. \textit{Distratores:} (A) ignora $V=1$ no Setor III; (B) inverte $9$ e $6$; (D) erra a ordem entre II, III e IV; (E) classifica $R = 4$ como Baixo (a faixa Baixo vai de $1$ a $3$).
  \item \textbf{Gabarito: (C).} Como a área possui resiliência muito baixa e alteração topográfica irreversível, sendo destinada a um novo uso produtivo e seguro (viveiro e captação pluvial), aplica-se a \textbf{reabilitação ambiental} (modalidade de recuperação). \textit{Distratores:} (A) restauração exigiria aproximação da condição original; (B) não se trata de poluição; (D) e (E) são conceitos incompatíveis com o caso.
\end{enumerate}
\end{boxGabarito}

\section*{Parte IV: Espelho Dissertativo e Rubrica de Correção (Questões 9 e 10)}

\begin{boxGabarito}[Espelho Dissertativo: Questão 9 (Valor: 2,0 pontos)]
\textbf{Eixos obrigatórios e pontuação:}
\begin{itemize}[leftmargin=1.5em, itemsep=2pt]
  \item \textbf{Item (a), Ruptura sociedade--natureza (0,6 pt):} explicitar que a crise local não é mera falha técnica ou destino climático, mas reflexo da racionalidade econômica e científica moderna que separou sociedade e natureza (Leff, 2001; Porto-Gonçalves, 2006), tratando água e vegetação como insumos apropriáveis de forma desigual.
  \item \textbf{Item (b), Geografia Socioambiental (0,6 pt):} demonstrar que a análise puramente naturalista oculta os conflitos pelo uso da água e as decisões econômicas, enquanto a análise exclusivamente social omite o funcionamento hidrológico e a degradação das cabeceiras; a Geografia Socioambiental (Mendonça, 2001) integra ambas as dimensões como unidade inseparável.
  \item \textbf{Item (c), Aplicação do Sistema GTP (0,8 pt):} organizar claramente os três polos de Bertrand (2004) e suas interações: \textbf{G (Geossistema)}, balanço hídrico, infiltração, desmatamento de cabeceiras e erosão; \textbf{T (Território)}, outorga/disputa pelo uso da água, estrutura fundiária, normas de APP e gestão da bacia; \textbf{P (Paisagem)}, transformação fisionômica do vale e perda de memória/identidade cultural dos moradores com o riacho.
\end{itemize}
\end{boxGabarito}

\begin{boxGabarito}[Espelho Dissertativo: Questão 10 (Valor: 2,0 pontos)]
\textbf{Eixos obrigatórios e pontuação:}
\begin{itemize}[leftmargin=1.5em, itemsep=2pt]
  \item \textbf{Item (a), Cadeia causal (0,5 pt):} \textit{Atividade:} abertura de estrada de acesso em vertente declivosa; \textit{Aspecto ambiental:} remoção da vegetação e exposição do solo sem drenagem adequada; \textit{Impacto direto:} erosão acelerada / ravinamento na encosta; \textit{Impacto indireto:} carreamento de sedimentos e assoreamento do açude comunitário a jusante (com redução do volume útil de água).
  \item \textbf{Item (b), Degradação vs. Poluição (0,4 pt):} configura inequivocamente \textbf{degradação ambiental} (alteração adversa das propriedades do solo e do corpo hídrico, PNMA Art. 3º, II). Caso o aporte de sedimentos eleve a turbidez e os sólidos suspensos acima dos padrões legais, prejudicando o abastecimento humano, a biota aquática e as condições sanitárias da água (PNMA Art. 3º, III), também se enquadra como \textbf{poluição hídrica} (por poluição física/sedimentar).
  \item \textbf{Item (c), Cálculos de Risco e $\Delta R\%$ (0,7 pt):}
  \begin{itemize}[leftmargin=1.3em, itemsep=1pt]
    \item Cenário 1: $R_1 = P_1 \times V_1 = 4 \times 3 = \mathbf{12}$ (\textbf{Risco Crítico}).
    \item Cenário 2: $R_2 = P_2 \times V_2 = 2 \times 2 = \mathbf{4}$ (\textbf{Risco Médio}).
    \item Variação percentual: $\Delta R\% = \frac{4 - 12}{12} \times 100\% = \frac{-8}{12} \times 100\% \approx \mathbf{-66,67\%}$ (redução de $66,67\%$ no índice de risco).
  \end{itemize}
  \item \textbf{Item (d), Restauração vs. Reabilitação na APP do açude (0,4 pt):} como se trata de \textbf{Área de Preservação Permanente (APP)} com função ecológica essencial de retenção de sedimentos e proteção hídrica, a diretriz prioritária é a \textbf{restauração ecológica} da faixa ciliar com espécies nativas da Caatinga, aproximando a estrutura e a função protetora do ecossistema original (admitindo-se ações estruturais auxiliares de bioengenharia para conter os sedimentos).
\end{itemize}
\end{boxGabarito}

\section*{Parte V: Tabela de Diagnóstico de Revisão do Módulo 1}

Marque abaixo o seu desempenho em cada questão. Se você errou ou teve dúvida em algum item, retorne imediatamente ao capítulo indicado antes de iniciar o Módulo 2:

\begin{center}
\renewcommand{\arraystretch}{1.22}
\small
\begin{tabularx}{\textwidth}{@{}c c p{3.8cm} >{\raggedright\arraybackslash}X c@{}}
  \toprule
  \rowcolor{uefsNavy}
  \textcolor{white}{\bfseries Questão} & \textcolor{white}{\bfseries Revisar} & \textcolor{white}{\bfseries Tema Central} & \textcolor{white}{\bfseries Referência-Base e Ação Recomendada} & \textcolor{white}{\bfseries Acertou?} \\
  \midrule
  \textbf{Q1} & \textbf{Cap. 1} & Crise ambiental e racionalidade & Leff (2001, cap. 2). Reler Etapas 3 e 9 do Cap. 1. & [ \quad ] \\
  \rowcolor{bgGray}
  \textbf{Q2} & \textbf{Cap. 1--2} & Modernidade e natureza social & Porto-Gonçalves (2006). Revisar transição sujeito--objeto (Cap. 2). & [ \quad ] \\
  \textbf{Q3} & \textbf{Cap. 2} & Conflitos socioambientais & Porto-Gonçalves (2006); Mendonça (2001). Revisar casos do Cap. 2 (EO1). & [ \quad ] \\
  \rowcolor{bgGray}
  \textbf{Q4} & \textbf{Cap. 3} & Geografia Socioambiental & Mendonça (2001). Reler Seção 3.2 e Erros 1--2 do Cap. 3 (EO2). & [ \quad ] \\
  \textbf{Q5} & \textbf{Cap. 3} & Sistema Tripolar GTP & Bertrand (2004). Revisar Figura~\ref{fig:triade-gtp} e Exemplos 3.2--3.3. & [ \quad ] \\
  \rowcolor{bgGray}
  \textbf{Q6} & \textbf{Cap. 4} & Aspecto, impacto e sinergia & Sánchez (2013, cap. 1); CONAMA 01/86. Reler Seção 4.2 e Erro 1 (EO3). & [ \quad ] \\
  \textbf{Q7} & \textbf{Cap. 4} & Matriz de Risco ($R = P \times V$) & Sánchez (2013). Refazer cálculo da Figura~\ref{fig:aspecto-risco} e Exemplo 4.3. & [ \quad ] \\
  \rowcolor{bgGray}
  \textbf{Q8} & \textbf{Cap. 4} & Reabilitação vs. Restauração & Sánchez (2013); Lei do SNUC. Reler Seção 4.5, Exemplo 4.5 e Erro 4. & [ \quad ] \\
  \textbf{Q9} & \textbf{Caps. 1--3} & Síntese dissertativa EO1--EO2 & Articular Leff, Porto-Gonçalves, Mendonça e Bertrand em parágrafos integrados. & [ \quad ] \\
  \rowcolor{bgGray}
  \textbf{Q10} & \textbf{Caps. 3--4} & Diagnóstico técnico EO3 & Treinar cadeia causal, PNMA, matriz $R = P \times V$ e escolha de PRAD. & [ \quad ] \\
  \bottomrule
\end{tabularx}
\end{center}
